% Auto-generated by extract-marking.py -- do not edit by hand unless you know what to do.
% Works for BOTH \documentclass[...,type=solutions,...] and
% [...,type=marking,...] compiles, chosen at compile time via
% \OlympMarkPick / \OlympMarkBeginTable (see olympiad-marking.sty).
% \eqref{n} falls back to a plain '(n)' whenever the referenced
% equation isn't actually typeset in the current document.

\OlympMarkBeginTable{|>{\centering\arraybackslash}p{0.065\linewidth}|>{\arraybackslash}p{0.74\linewidth}|>{\centering\arraybackslash}p{0.05\linewidth}|>{\centering\arraybackslash}p{0.045\linewidth}|}{|>{\centering\arraybackslash}p{0.055\linewidth}|>{\arraybackslash}p{0.54\linewidth}|>{\centering\arraybackslash}p{0.04\linewidth}|>{\centering\arraybackslash}p{0.045\linewidth}|>{\centering\arraybackslash}p{0.04\linewidth}|>{\centering\arraybackslash}p{0.04\linewidth}|>{\centering\arraybackslash}p{0.04\linewidth}|}
\hline
\OlympMarkPick{\textbf{\textnumero} & \multicolumn{1}{c|}{\textbf{Содержание}} & \multicolumn{2}{c|}{\textbf{Баллы}} \\ \hline}{\textbf{\textnumero} & \multicolumn{1}{c|}{\textbf{Содержание}} & \multicolumn{2}{c|}{\textbf{Баллы}} & \multicolumn{2}{c|}{\textbf{MP}} & \textbf{A}\\ \hline}
\OlympMarkPick{\multirow{4}{*}{\textbf{\OlympCurrentPrefix.1}} & Формула \eqref{1}: $H\cdot l = NI,$ & 0.2 & \multirow{4}{*}{\textbf{0.4}} \\ \cline{2-3}}{\multirow{4}{*}{\textbf{\OlympCurrentPrefix.1}} & Формула \eqref{1}: $H\cdot l = NI,$ & 0.2 & \multirow{4}{*}{\textbf{0.4}} & & \multirow{4}{*}{} & \multirow{4}{*}{} \\ \cline{2-3}\cline{5-5}}
\OlympMarkPick{ & Формула \eqref{2}: $B = \frac{\mu\mu_0 NI}{l}=\qty{1}{\T}.$ & 0.1 &  \\ \cline{2-3}}{ & Формула \eqref{2}: $B = \frac{\mu\mu_0 NI}{l}=\qty{1}{\T}.$ & 0.1 &  & & &  \\ \cline{2-3}\cline{5-5}}
\OlympMarkPick{ & Численный ответ в \eqref{2} & 0.1 &  \\ \cline{2-3}}{ & Численный ответ в \eqref{2} & 0.1 &  & & &  \\ \cline{2-3}\cline{5-5}}
\OlympMarkPick{ & \quad\textit{Если участник сразу записал \eqref{1} через $B=\mu\mu_0 H$, даётся балл за \eqref{1} тоже} &  &  \\ \hline}{ & \quad\textit{Если участник сразу записал \eqref{1} через $B=\mu\mu_0 H$, даётся балл за \eqref{1} тоже} &  &  & & &  \\ \hline}
\OlympMarkPick{\multirow{3}{*}{\textbf{\OlympCurrentPrefix.2}} & Формула \eqref{3}: $\Phi = B\cdot \pi r^2.$ & 0.2 & \multirow{3}{*}{\textbf{0.6}} \\ \cline{2-3}}{\multirow{3}{*}{\textbf{\OlympCurrentPrefix.2}} & Формула \eqref{3}: $\Phi = B\cdot \pi r^2.$ & 0.2 & \multirow{3}{*}{\textbf{0.6}} & & \multirow{3}{*}{} & \multirow{3}{*}{} \\ \cline{2-3}\cline{5-5}}
\OlympMarkPick{ & Формула \eqref{4}: $R_m = \frac l{\mu\mu_0\cdot\pi r^2}=\qty{6.33e5}{\A\per\Wb}.$ & 0.2 &  \\ \cline{2-3}}{ & Формула \eqref{4}: $R_m = \frac l{\mu\mu_0\cdot\pi r^2}=\qty{6.33e5}{\A\per\Wb}.$ & 0.2 &  & & &  \\ \cline{2-3}\cline{5-5}}
\OlympMarkPick{ & Численный ответ в \eqref{4} & 0.2 &  \\ \hline}{ & Численный ответ в \eqref{4} & 0.2 &  & & &  \\ \hline}
\OlympMarkPick{\multirow{3}{*}{\textbf{\OlympCurrentPrefix.3}} & Формула \eqref{5}: $\Psi=N\Phi,$ & 0.2 & \multirow{3}{*}{\textbf{0.6}} \\ \cline{2-3}}{\multirow{3}{*}{\textbf{\OlympCurrentPrefix.3}} & Формула \eqref{5}: $\Psi=N\Phi,$ & 0.2 & \multirow{3}{*}{\textbf{0.6}} & & \multirow{3}{*}{} & \multirow{3}{*}{} \\ \cline{2-3}\cline{5-5}}
\OlympMarkPick{ & Формула \eqref{6}: $L = \frac{N^2}{R_m} = \frac{\mu\mu_0\pi r^2N^2}l =\qty{63.2}{\mH}.$ & 0.2 &  \\ \cline{2-3}}{ & Формула \eqref{6}: $L = \frac{N^2}{R_m} = \frac{\mu\mu_0\pi r^2N^2}l =\qty{63.2}{\mH}.$ & 0.2 &  & & &  \\ \cline{2-3}\cline{5-5}}
\OlympMarkPick{ & Численный ответ в \eqref{6} & 0.2 &  \\ \hline}{ & Численный ответ в \eqref{6} & 0.2 &  & & &  \\ \hline}
\OlympMarkPick{\multirow{6}{*}{\textbf{\OlympCurrentPrefix.4}} & Формула \eqref{7}: $B_0\cos\alpha=B\cos\beta,$ & 0.3 & \multirow{6}{*}{\textbf{1.2}} \\ \cline{2-3}}{\multirow{6}{*}{\textbf{\OlympCurrentPrefix.4}} & Формула \eqref{7}: $B_0\cos\alpha=B\cos\beta,$ & 0.3 & \multirow{6}{*}{\textbf{1.2}} & & \multirow{6}{*}{} & \multirow{6}{*}{} \\ \cline{2-3}\cline{5-5}}
\OlympMarkPick{ & Формула \eqref{8}: $B_0\sin\alpha = \frac B\mu \sin\beta.$ & 0.3 &  \\ \cline{2-3}}{ & Формула \eqref{8}: $B_0\sin\alpha = \frac B\mu \sin\beta.$ & 0.3 &  & & &  \\ \cline{2-3}\cline{5-5}}
\OlympMarkPick{ & Формула \eqref{9}: $\tan\beta = \mu \tan\alpha,$ & 0.2 &  \\ \cline{2-3}}{ & Формула \eqref{9}: $\tan\beta = \mu \tan\alpha,$ & 0.2 &  & & &  \\ \cline{2-3}\cline{5-5}}
\OlympMarkPick{ & Численное значение \eqref{10}: $\beta \approx \ang{89.83},$ & 0.1 &  \\ \cline{2-3}}{ & Численное значение \eqref{10}: $\beta \approx \ang{89.83},$ & 0.1 &  & & &  \\ \cline{2-3}\cline{5-5}}
\OlympMarkPick{ & Формула \eqref{11}: $\frac B{B_0}=\sqrt{\cos^2\alpha+\mu^2\sin^2\alpha} \approx 347.$ & 0.2 &  \\ \cline{2-3}}{ & Формула \eqref{11}: $\frac B{B_0}=\sqrt{\cos^2\alpha+\mu^2\sin^2\alpha} \approx 347.$ & 0.2 &  & & &  \\ \cline{2-3}\cline{5-5}}
\OlympMarkPick{ & Численный ответ в \eqref{11} & 0.1 &  \\ \hline}{ & Численный ответ в \eqref{11} & 0.1 &  & & &  \\ \hline}
\OlympMarkPick{\multirow{5}{*}{\textbf{\OlympCurrentPrefix.5}} & Формула \eqref{12}: $H_1\cdot(4a-d) + H_2\cdot d = NI,$ & 0.2 & \multirow{5}{*}{\textbf{1.0}} \\ \cline{2-3}}{\multirow{5}{*}{\textbf{\OlympCurrentPrefix.5}} & Формула \eqref{12}: $H_1\cdot(4a-d) + H_2\cdot d = NI,$ & 0.2 & \multirow{5}{*}{\textbf{1.0}} & & \multirow{5}{*}{} & \multirow{5}{*}{} \\ \cline{2-3}\cline{5-5}}
\OlympMarkPick{ & Формула \eqref{13}: $B = \mu_0 H_2 = \frac{\mu_0}\mu H_1.$ & 0.3 &  \\ \cline{2-3}}{ & Формула \eqref{13}: $B = \mu_0 H_2 = \frac{\mu_0}\mu H_1.$ & 0.3 &  & & &  \\ \cline{2-3}\cline{5-5}}
\OlympMarkPick{ & Формула \eqref{14}: $\Phi = \frac{\mu_0 NI S}{d + \mu(4a-d)}.$ & 0.3 &  \\ \cline{2-3}}{ & Формула \eqref{14}: $\Phi = \frac{\mu_0 NI S}{d + \mu(4a-d)}.$ & 0.3 &  & & &  \\ \cline{2-3}\cline{5-5}}
\OlympMarkPick{ & Формула \eqref{15}: $R_{m,1} = \frac{4a-d}{\mu\mu_0S},$ & 0.1 &  \\ \cline{2-3}}{ & Формула \eqref{15}: $R_{m,1} = \frac{4a-d}{\mu\mu_0S},$ & 0.1 &  & & &  \\ \cline{2-3}\cline{5-5}}
\OlympMarkPick{ & Формула \eqref{16}: $R_{m,2} = \frac d{\mu_0S},$ & 0.1 &  \\ \hline}{ & Формула \eqref{16}: $R_{m,2} = \frac d{\mu_0S},$ & 0.1 &  & & &  \\ \hline}
\OlympMarkPick{\multirow{2}{*}{\textbf{\OlympCurrentPrefix.6}} & Формула \eqref{17}: $d = \frac{4a}{\mu+1}\approx\frac{4a}{\mu} = \qty{0.4}{\mm}.$ & 0.2 & \multirow{2}{*}{\textbf{0.4}} \\ \cline{2-3}}{\multirow{2}{*}{\textbf{\OlympCurrentPrefix.6}} & Формула \eqref{17}: $d = \frac{4a}{\mu+1}\approx\frac{4a}{\mu} = \qty{0.4}{\mm}.$ & 0.2 & \multirow{2}{*}{\textbf{0.4}} & & \multirow{2}{*}{} & \multirow{2}{*}{} \\ \cline{2-3}\cline{5-5}}
\OlympMarkPick{ & Численный ответ в \eqref{17} & 0.2 &  \\ \hline}{ & Численный ответ в \eqref{17} & 0.2 &  & & &  \\ \hline}
\OlympMarkPick{\multirow{8}{*}{\textbf{\OlympCurrentPrefix.7}} & В пунктах \solref{7}, \solref{10}--\solref{11}, связанные с законами магнитных цепей критерии (\eqref{18}, \eqref{20}, \eqref{27}--\eqref{29}, \eqref{33}--\eqref{34}) принимаются при наличии аналогичных выражений & \textbf{?} & \multirow{8}{*}{\textbf{1.5}} \\ \cline{2-3}}{\multirow{8}{*}{\textbf{\OlympCurrentPrefix.7}} & В пунктах \solref{7}, \solref{10}--\solref{11}, связанные с законами магнитных цепей критерии (\eqref{18}, \eqref{20}, \eqref{27}--\eqref{29}, \eqref{33}--\eqref{34}) принимаются при наличии аналогичных выражений & \textbf{?} & \multirow{8}{*}{\textbf{1.5}} & & \multirow{8}{*}{} & \multirow{8}{*}{} \\ \cline{2-3}\cline{5-5}}
\OlympMarkPick{ & Формула \eqref{18}: $\Phi(t) = \frac{N_1 I_1(t)}{R_m},$ & 0.2 &  \\ \cline{2-3}}{ & Формула \eqref{18}: $\Phi(t) = \frac{N_1 I_1(t)}{R_m},$ & 0.2 &  & & &  \\ \cline{2-3}\cline{5-5}}
\OlympMarkPick{ & Формула \eqref{19}: $\Psi_2(t) = N_2\Phi(t),$ & 0.2 &  \\ \cline{2-3}}{ & Формула \eqref{19}: $\Psi_2(t) = N_2\Phi(t),$ & 0.2 &  & & &  \\ \cline{2-3}\cline{5-5}}
\OlympMarkPick{ & Формула \eqref{20}: $M = \frac{N_1N_2}{R_m} = \frac{\mu\mu_0N_1N_2S}{4a}.$ & 0.2 &  \\ \cline{2-3}}{ & Формула \eqref{20}: $M = \frac{N_1N_2}{R_m} = \frac{\mu\mu_0N_1N_2S}{4a}.$ & 0.2 &  & & &  \\ \cline{2-3}\cline{5-5}}
\OlympMarkPick{ & Формула \eqref{21}: $V(t) = -N_1\frac{d\Phi}{dt},$ & 0.2 &  \\ \cline{2-3}}{ & Формула \eqref{21}: $V(t) = -N_1\frac{d\Phi}{dt},$ & 0.2 &  & & &  \\ \cline{2-3}\cline{5-5}}
\OlympMarkPick{ & Формула \eqref{22}: $I_2 R = \frac{N_1}{N_2}V(t).$ & 0.2 &  \\ \cline{2-3}}{ & Формула \eqref{22}: $I_2 R = \frac{N_1}{N_2}V(t).$ & 0.2 &  & & &  \\ \cline{2-3}\cline{5-5}}
\OlympMarkPick{ & Формула \eqref{23}: $N_1I_1+N_2I_2 = 0.$ & 0.2 &  \\ \cline{2-3}}{ & Формула \eqref{23}: $N_1I_1+N_2I_2 = 0.$ & 0.2 &  & & &  \\ \cline{2-3}\cline{5-5}}
\OlympMarkPick{ & Формула \eqref{24}: $R'=R\ab(\frac{N_1}{N_2})^2$ & 0.3 &  \\ \hline}{ & Формула \eqref{24}: $R'=R\ab(\frac{N_1}{N_2})^2$ & 0.3 &  & & &  \\ \hline}
\OlympMarkPick{}{\pagebreak}
\OlympMarkPick{\multirow{3}{*}{\textbf{\OlympCurrentPrefix.8}} & Записано: <<Сила тока $I$>> & 0.1 & \multirow{3}{*}{\textbf{0.5}} \\ \cline{2-3}}{\multirow{3}{*}{\textbf{\OlympCurrentPrefix.8}} & Записано: <<Сила тока $I$>> & 0.1 & \multirow{3}{*}{\textbf{0.5}} & & \multirow{3}{*}{} & \multirow{3}{*}{} \\ \cline{2-3}\cline{5-5}}
\OlympMarkPick{ & Записано: <<Напряжённость магнитного поля $\vec H$>> & 0.2 &  \\ \cline{2-3}}{ & Записано: <<Напряжённость магнитного поля $\vec H$>> & 0.2 &  & & &  \\ \cline{2-3}\cline{5-5}}
\OlympMarkPick{ & Записано: <<Удельная проводимость $\sigma$ или $1/\rho$>> & 0.2 &  \\ \hline}{ & Записано: <<Удельная проводимость $\sigma$ или $1/\rho$>> & 0.2 &  & & &  \\ \hline}
\OlympMarkPick{\multirow{2}{*}{\textbf{\OlympCurrentPrefix.9}} & Формула \eqref{25}: $w_m = \frac{BH}{2} = \frac{B^2}{2\mu\mu_0}=\frac{\mu\mu_0 H^2}{2}.$ & 0.2 & \multirow{2}{*}{\textbf{0.5}} \\ \cline{2-3}}{\multirow{2}{*}{\textbf{\OlympCurrentPrefix.9}} & Формула \eqref{25}: $w_m = \frac{BH}{2} = \frac{B^2}{2\mu\mu_0}=\frac{\mu\mu_0 H^2}{2}.$ & 0.2 & \multirow{2}{*}{\textbf{0.5}} & & \multirow{2}{*}{} & \multirow{2}{*}{} \\ \cline{2-3}\cline{5-5}}
\OlympMarkPick{ & Формула \eqref{26}: $W_M = \frac{\mathcal F^2}{2R_m}$ & 0.3 &  \\ \hline}{ & Формула \eqref{26}: $W_M = \frac{\mathcal F^2}{2R_m}$ & 0.3 &  & & &  \\ \hline}
\OlympMarkPick{\multirow{6}{*}{\textbf{\OlympCurrentPrefix.10}} & Формула \eqref{27}: $NI_1(t) + NI_2(t) = 3R_m(\Phi_1+\Phi_2)$ & 0.3 & \multirow{6}{*}{\textbf{1.5}} \\ \cline{2-3}}{\multirow{6}{*}{\textbf{\OlympCurrentPrefix.10}} & Формула \eqref{27}: $NI_1(t) + NI_2(t) = 3R_m(\Phi_1+\Phi_2)$ & 0.3 & \multirow{6}{*}{\textbf{1.5}} & & \multirow{6}{*}{} & \multirow{6}{*}{} \\ \cline{2-3}\cline{5-5}}
\OlympMarkPick{ & Формула \eqref{28}: $NI_1(t)+N_0I_M(t) = 3R_m\cdot\Phi_1 + R_m\cdot (\Phi_1-\Phi_2).$ & 0.3 &  \\ \cline{2-3}}{ & Формула \eqref{28}: $NI_1(t)+N_0I_M(t) = 3R_m\cdot\Phi_1 + R_m\cdot (\Phi_1-\Phi_2).$ & 0.3 &  & & &  \\ \cline{2-3}\cline{5-5}}
\OlympMarkPick{ & Формула \eqref{29}: $\Delta\Phi = \frac{N(I_1(t)-I_2(t))}{5R_m}.$ & 0.3 &  \\ \cline{2-3}}{ & Формула \eqref{29}: $\Delta\Phi = \frac{N(I_1(t)-I_2(t))}{5R_m}.$ & 0.3 &  & & &  \\ \cline{2-3}\cline{5-5}}
\OlympMarkPick{ & Формула \eqref{30}: $I_M(t) = -\frac{N_0}R\frac{d}{dt}\Delta\Phi,$ & 0.1 &  \\ \cline{2-3}}{ & Формула \eqref{30}: $I_M(t) = -\frac{N_0}R\frac{d}{dt}\Delta\Phi,$ & 0.1 &  & & &  \\ \cline{2-3}\cline{5-5}}
\OlympMarkPick{ & Формула \eqref{31}: $I_M = \frac{2\pi fNN_0(I_1-I_2)}{\sqrt{(5R\rho_m l)^2+(4\pi fN_0^2)^2}}=\qty{11.8}{\mA}.$ & 0.3 &  \\ \cline{2-3}}{ & Формула \eqref{31}: $I_M = \frac{2\pi fNN_0(I_1-I_2)}{\sqrt{(5R\rho_m l)^2+(4\pi fN_0^2)^2}}=\qty{11.8}{\mA}.$ & 0.3 &  & & &  \\ \cline{2-3}\cline{5-5}}
\OlympMarkPick{ & Численный ответ в \eqref{31} & 0.2 &  \\ \hline}{ & Численный ответ в \eqref{31} & 0.2 &  & & &  \\ \hline}
\OlympMarkPick{\multirow{9}{*}{\textbf{\OlympCurrentPrefix.11}} & Формула \eqref{32}: $V=\pi a^2(d-x).$ & 0.1 & \multirow{9}{*}{\textbf{1.8}} \\ \cline{2-3}}{\multirow{9}{*}{\textbf{\OlympCurrentPrefix.11}} & Формула \eqref{32}: $V=\pi a^2(d-x).$ & 0.1 & \multirow{9}{*}{\textbf{1.8}} & & \multirow{9}{*}{} & \multirow{9}{*}{} \\ \cline{2-3}\cline{5-5}}
\OlympMarkPick{ & Формула \eqref{33}: $W_m = \frac{\mathcal F^2}{2R_m},$ & 0.1 &  \\ \cline{2-3}}{ & Формула \eqref{33}: $W_m = \frac{\mathcal F^2}{2R_m},$ & 0.1 &  & & &  \\ \cline{2-3}\cline{5-5}}
\OlympMarkPick{ & Формула \eqref{34}: $R_m = \frac{d-x}{\mu_0 \pi a^2}.$ & 0.4 &  \\ \cline{2-3}}{ & Формула \eqref{34}: $R_m = \frac{d-x}{\mu_0 \pi a^2}.$ & 0.4 &  & & &  \\ \cline{2-3}\cline{5-5}}
\OlympMarkPick{ & \quad\textit{Участник может записать вклады в энергию от сердечников или боковых зазоров; балл ставится если в их выражении имеется правильное слагаемое \eqref{33}--\eqref{34}} &  &  \\ \cline{2-3}}{ & \quad\textit{Участник может записать вклады в энергию от сердечников или боковых зазоров; балл ставится если в их выражении имеется правильное слагаемое \eqref{33}--\eqref{34}} &  &  & & &  \\ \cline{2-3}\cline{5-5}}
\OlympMarkPick{ & Формула \eqref{35}: $F=-\frac{dW_m}{dx},$ & 0.2 &  \\ \cline{2-3}}{ & Формула \eqref{35}: $F=-\frac{dW_m}{dx},$ & 0.2 &  & & &  \\ \cline{2-3}\cline{5-5}}
\OlympMarkPick{ & Формула \eqref{36}: $F(x_c) \geq kx_c.$ & 0.2 &  \\ \cline{2-3}}{ & Формула \eqref{36}: $F(x_c) \geq kx_c.$ & 0.2 &  & & &  \\ \cline{2-3}\cline{5-5}}
\OlympMarkPick{ & Формула \eqref{37}: $\mathcal F_0 = \frac{d-x_c}a\sqrt\frac{2kx_c}{\mu_0\pi}$ & 0.5 &  \\ \cline{2-3}}{ & Формула \eqref{37}: $\mathcal F_0 = \frac{d-x_c}a\sqrt\frac{2kx_c}{\mu_0\pi}$ & 0.5 &  & & &  \\ \cline{2-3}\cline{5-5}}
\OlympMarkPick{ & Численное значение \eqref{38}: $I_M = \qty{11.25}{\mA}.$ & 0.2 &  \\ \cline{2-3}}{ & Численное значение \eqref{38}: $I_M = \qty{11.25}{\mA}.$ & 0.2 &  & & &  \\ \cline{2-3}\cline{5-5}}
\OlympMarkPick{ & Показано, что $I_M<\qty{11.8}{\mA}$ & 0.1 &  \\ \hline}{ & Показано, что $I_M<\qty{11.8}{\mA}$ & 0.1 &  & & &  \\ \hline}
\OlympMarkPick{\textbf{Итого} &  &  & \textbf{10.0} \\ \hline}{\textbf{Итого} &  &  & \textbf{10.0} &  &  &  \\ \hline}
\end{tabularx}
